% !TEX program = lualatex
%% ============================================================
%% MODELO DE APOSTILA (material didático da NeuroDynamics)
%%
%% 1. Peça o código ao controle de documentos (NRO-PUB-001) e preencha
%%    \documento{...} abaixo. O código sai como NRO-XXX-000-N.
%% 2. Compile com LuaLaTeX (no Overleaf: Menu, Compilador, LuaLaTeX; na
%%    linha de comando: latexmk -lualatex modelo-apostila.tex). Compile
%%    duas vezes: o sumário, os totais e a ficha de progresso vêm da
%%    compilação anterior.
%% 3. Troque o que está entre colchetes e apague o que não usar. As
%%    partes, os níveis e os tópicos se numeram sozinhos.
%% O exemplo completo é o NRO-PRO-014-1.tex.
%% ============================================================
\documentclass{nro-apostila}

\documento{
  codigo        = NRO-XXX-000,            % do NRO-PUB-001
  numero        = 1,                      % a cópia, no fim do código
  revisao       = A,                      % A, B, C... a cada alteração
  emissor       = {Departamento de Pesquisa e Desenvolvimento},
  tipo          = {Manual ou guia},       % o subtipo do NRO-PUB-002
  categoria     = {Material didático},    % o rótulo da capa
  titulo        = {Título da apostila},
  subtitulo     = {Uma frase dizendo para quem é e o que ensina.},
  eixos         = {Área um\hspace{9mm}Área dois\hspace{9mm}Área três},
  data          = {27 de outubro de 2026},
  redator       = {Nome de quem escreveu},
  revisor       = {},
  aprovador     = {},
  rotuloproduto = {No produto},           % o rótulo do quadro de aplicação
  autoria       = {NeuroDynamics},
}

\begin{document}

\capa
\sumario

%% ---------- páginas de introdução ----------
\pagina{Introdução}{Boas-vindas}
\lede{Um parágrafo de abertura, maior, que diz o que o leitor vai encontrar.}

Texto corrido. Use \textbf{negrito} só para termos que o leitor precisa guardar, \code{codigo\_inline} para nomes de arquivo, comandos e registradores, e \cod{NRO-XXX-000} para códigos de documento (nunca quebram na linha).

\begin{estatisticas}[4]
  \estatistica{\totalpartes}{Partes}
  \estatistica{\totaltopicos}{Tópicos}
  \estatistica{3}{Níveis}
  \estatistica{\totalexercicios}{Exercícios}
\end{estatisticas}

\secao{Princípios}
\begin{cartoes}[2]
  \cartao{Primeiro princípio}{Uma frase.}
  \cartao{Segundo princípio}{Uma frase.}
\end{cartoes}

\begin{nota}
  Uma nota neutra, para um esclarecimento.
\end{nota}
\begin{atencao}
  Um alerta, para o que não pode ser ignorado (segurança, risco, dado perdido).
\end{atencao}
\begin{definicao}{Termo}
  A definição curta de um termo, em uma ou duas frases.
\end{definicao}

\secao{Uma tabela}
\begin{tabularx}{\linewidth}{@{}L{30mm}Y L{24mm}@{}}\reguatopo
  \cab{Coluna} & \cab{Outra coluna} & \cab{Código} \\ \reguameio
  Linha um  & Texto da célula & \cod{1.1} \\ \reguafina
  Linha dois & Texto da célula & \cod{1.2} \\ \reguabase
\end{tabularx}

%% ---------- uma parte: \iniciarpartes{1} antes do \capa começa na Parte 1 ----------
\begin{parte}{Título da parte}{Uma frase de abertura: o que esta parte cobre e para quem é.}
  \item O que o leitor sabe fazer ao terminar (um objetivo por item).
  \item Outro objetivo.
\end{parte}

%% um nível (opcional): a capa da parte mostra os níveis sozinha; sem nível, mostra a lista dos tópicos
\nivel{Básico}{O que este nível cobre.}

\topico{Título do tópico}
O conceito explicado de forma geral, válido para qualquer projeto, em duas ou três frases.

\begin{noproduto}
  Onde o conceito aparece no produto: repositório, esquemático, módulo.
\end{noproduto}

\begin{paraestudar}
  \item \leitura{https://exemplo.org}{Título da fonte}
  \item Livro: Título (autor), capítulos tais
\end{paraestudar}

\begin{pratique}
  \item Exercício conceitual.
  \item Exercício prático, ligado ao código ou ao hardware real.
\end{pratique}

\topico{Outro tópico}
Texto.

%% ---------- apêndices ----------
\apendice{Glossário}
\begin{glossario}
  \termo{Termo}{Significado em uma frase.}
\end{glossario}

\apendice{Ficha de progresso}
Marque cada tópico concluído e anote a data.
\par\vspace{2mm}
\fichaprogresso

\apendice{Controle do documento}
\controledocumento
\secao{Histórico de revisões}
\begin{historico}
  \linharevisao{27/10/2026}{A}{Nome}{Primeira edição.}
\end{historico}

\contracapa{Uma frase de fechamento: onde reportar erros e como contribuir.}

\end{document}
